% ===========================================================================
\section{The Fourier basis}\label{sec:basis}
\warmth{0}\quad\whisper{A signal is a vector; the pure tones are an orthonormal basis; the transform is the change of coordinates.}

\begin{strand}
Think of a signal $f$ as a vector in an (infinite-dimensional) inner-product space. The
\keyword{pure tones} $e_\xi(x)=e^{2\pi i\xi x}$ are a basis of that space, one vector per frequency
$\xi$. The Fourier transform is nothing but reading off coordinates in this basis: each
$\widehat f(\xi)$ is the inner product of $f$ with the tone $e_\xi$, i.e.\ ``how much of frequency
$\xi$ is in $f$''.
\end{strand}

\block{The transform}

\begin{definition}[Fourier transform]\label{def:ft}
For $f\colon\R\to\C$ (integrable enough),
\[
  \widehat f(\xi)\;=\;\F f(\xi)\;=\;\int_{-\infty}^{\infty} f(x)\,e^{\fillin[1.6cm]}\dd x ,
\]
the coordinate of $f$ along the tone $e_\xi$. The finite version is the \keyword{DFT}:
$\widehat x_k=\sum_{n=0}^{N-1} x_n\,e^{-2\pi i kn/N}$.
\end{definition}

\recall{Why complex exponentials and not $\sin,\cos$? Because $e_\xi$ are the \emph{eigenfunctions of shift and of $d/dx$}: shifting or differentiating only multiplies them by a scalar. That is the entire reason Fourier diagonalizes those operations.}

\block{Three transforms to know cold}
\begin{yourturn}
Match each signal to its transform (these recur everywhere):
\begin{itemize}[leftmargin=1.3em,itemsep=1pt]
  \item a single tone $f(x)=e^{2\pi i\xi_0 x}$ \;$\longmapsto$\; a spike at $\fillin[1.6cm]$;
  \item the box $\mathbf 1_{[-1/2,\,1/2]}(x)$ \;$\longmapsto$\; $\widehat f(\xi)=\fillin[2cm]$ (the $\sinc$);
  \item the Gaussian $e^{-\pi x^2}$ \;$\longmapsto$\; $\fillin[2cm]$ (itself!).
\end{itemize}
\recall{Answers: $\delta(\xi-\xi_0)$; $\ \sinc(\xi)=\sin(\pi\xi)/(\pi\xi)$; $\ e^{-\pi\xi^2}$. The Gaussian is the transform's fixed point, and narrow box $\leftrightarrow$ wide $\sinc$ is uncertainty in miniature (§\ref{sec:why}).}
\end{yourturn}

\begin{strand}
Two morals already visible. First, a perfectly localized signal in $x$ (the spike/tone) is
perfectly spread in $\xi$, and vice versa. Second, the transform is a \keyword{change of
orthonormal basis}, so it loses no information and preserves lengths; that is Parseval's theorem,
coming in §\ref{sec:dict}--\ref{sec:why}.
\end{strand}

\loose{Conventions vary: some put $2\pi$ in front, some in the exponent. We keep it in the exponent, which makes the convolution theorem and Parseval cleanest (no stray constants). Always check a book's convention before borrowing a formula.}
